\centerline{\bf \Huge Эйлеровы графы.}
\title{Эйлеровы графы.}

1. Из доски $4\times4$ вырезали угловые клетки. Может ли шахматный конь обойти все клетки доски, пройдя по каждому полю ровно по одному разу?

2. В Элисте 993 дороги. Очир может обойти всю Элисту так, что по каждой дороге он пройдет ровно два раза, но не может обойти так, чтобы пройти ровно по одному разу. Как такое может быть?

3. При каких $n$ правильный $n$-угольник со всеми его диагоналями можно нарисовать не отрывая карандаша?

4. Город Токио представляет собой связный граф без кратных ребер, где вершины это дома, а ребра это дороги их соединяющие. Путешественница Ази хочет пройтись вдоль каждой дороги, причем на каждой побывать ровно два раза. Докажите, что ей удасться это сделать.

$5^*.$ У Паука Германа есть паутина, представляющаяя из себя связный граф из $2n$ вершин, соединенные нитями(ребрами), причем каждая вершина нечетной степени. Ползти Герман может только по своим нитям.  Пройдя еще раз по своей нити он добавляет к ней еще один слой. Чтобы поймать муху Алину ему нужно добавить к каждой паутине $n-1$ слой. Докажите, что пауку Герману удасться это сделать каким бы связным графом ни была его паутина.